\origpage{95}{GRAPH MINORS. XI}
\stmt
(7.9) \textit{If $\Sigma$ is connected, $X, Y \subseteq \Sigma$ are simply-connected, and $X \cap Y \ne \emptyset$ is arc-wise connected and nice, and $X \cup Y \ne \Sigma$, then $X \cap Y$ is simply-connected.}

\textit{Proof.} Let $z \in \Sigma - (X \cup Y)$. Let $F \subseteq X \cap Y$ be an $O$-arc. Then $F \subseteq X$ and $X$ is simply-connected, and so by (7.2) $F$ bounds a disc $\Delta_1 \subseteq X$. Similarly $F$ bounds a disc $\Delta_2 \subseteq Y$. Now $z \notin \Delta_1, \Delta_2$ since $z \notin X, Y$, and yet there is at most one disc in $\Sigma$ bounded by $F$ not containing $z$, since $\Sigma$ is connected. Thus $\Delta_1 = \Delta_2$, and $\Delta_1 \subseteq X \cap Y$. We deduce from (7.3) that $X \cap Y$ is simply-connected. $\square$

\stmt
(7.10) \textit{If $X \subseteq \Sigma$ is arc-wise connected, $Y \subseteq \Sigma$ is closed, $bd(Y)$ is arc-wise connected, and $\emptyset \ne bd(Y) \subseteq X$, then $X \cap Y$ is arc-wise connected.}

\textit{Proof.} Since $bd(Y)$ is arc-wise connected and non-empty, and $bd(Y) \subseteq X \cap Y$ since $Y$ is closed, there is an arc-wise connected component $C$ of $X \cap Y$ with $bd(Y) \subseteq C$. We claim that $C = X \cap Y$. For let $x \in X \cap Y$, let $y \in bd(Y)$, and let $\phi \colon [0, 1] \to X$ be continuous with $\phi(0) = x$ and $\phi(1) = y$. Choose $\alpha$ with $0 \le \alpha \le 1$, maximum so that $\phi(\beta) \in Y$ for all $\beta$ with $0 \le \beta \le \alpha$ (this is possible since $Y$ is closed). Then $\phi(\alpha) \in bd(Y)$; for if $\alpha < 1$ this follows from the maximality of $\alpha$, and if $\alpha = 1$ it is true since $y \in bd(Y)$. Moreover $\phi(\beta) \in X \cap Y$ for $0 \le \beta \le \alpha$, and so $x$ and $\phi(\alpha)$ belong to the same arc-wise connected component of $X \cap Y$. Since $\phi(\alpha) \in bd(Y) \subseteq C$ it follows that $x \in C$. This proves our claim that $C = X \cap Y$, and so $X \cap Y$ is arc-wise connected. $\square$

\stmt
(7.11) \textit{If $X \subseteq \Sigma$ is closed and arc-wise connected and $Y_1, \dots, Y_r \subseteq \Sigma$ are closed discs with boundary in $X$, and $Z = X \cup Y_1 \cup \dots \cup Y_r$, and every continuous $\phi \colon S^1 \to X$ is null-homotopic in $Z$, then every continuous $\phi \colon S^1 \to Z$ is null-homotopic in $Z$.}

\textit{Proof.} For $0 \le i \le r$ let $Z_i = X \cup Y_1 \cup \dots \cup Y_i$. We prove by induction on $i$ that $Z_i$ is arc-wise connected and every continuous $\phi \colon S^1 \to Z_i$ is null-homotopic in $Z$. Certainly this is true if $i = 0$. Suppose $1 \le i \le r$, and it holds for $i - 1$. By (7.10), since $Z_{i-1}$ is arc-wise connected, $Y_i$ is closed, $bd(Y_i)$ is arc-wise connected and $\emptyset \ne bd(Y_i) \subseteq Z_{i-1}$, it follows that $Z_{i-1} \cap Y_i$ is arc-wise connected. By (3.1), since $Y_i$ is simply-connected and $Z_{i-1}$, $Y_i$ are closed and $Z_{i-1} \cap Y_i \ne \emptyset$ is arc-wise connected, it follows that every continuous $\phi \colon S^1 \to Y_i \cup Z_{i-1}$ is homotopic in $Y_i \cup Z_{i-1}$ and hence in $Z$ to some $\phi \colon S^1 \to Z_{i-1}$, and hence is null-homotopic in $Z$. This completes the inductive argument, and the result follows. $\square$
